% Appendix A: Changes from Earlier Versions. Printed after the references when \ifappendix is true (main.tex). The
% build main_noappendix.tex leaves it out.
\section{Changes from Earlier Versions}\label{app:changes}
% arXiv versions 1 (28 September 2026) and 2 (29 September 2026) have the same source files and the same PDF text,
% apart from the arXiv stamp. Version 2 changed only the abstract of the arXiv listing, where "this http URL"
% (jax.vmap, jax.pmap) became "JAX's vmap" and "JAX's pmap".
This paper is version 3 of arXiv:2609.34818. Version 2 changed only how the abstract on the arXiv page wrote two JAX
function names. The paper itself is the same in versions 1 and 2, so every claim listed here appears in both. % src: arXiv listings of 2609.34818 as read on 2026-11-23 ("[v1] Mon, 28 Sep 2026", "[v2] Tue, 29 Sep 2026"), same e-print source files in v1 and v2
% The title of version 1, as in paper/preprint/source/main.tex (pdftitle).
Version 1 had the title \vonequote{Accelerator Choice Is Not Enough: AlphaFold2 Inference on Cloud TPUs.} Version 3
has a new title and a new structure.
We ran no new experiment for this version. We checked the claims and numbers of
version 1 against the committed result files. This appendix lists the claims that changed and the reason for each
change.

\subsection*{Ratios Between Platforms}
Version 1 reported ratios of 451 between one TPU chip and the CPU runtime and of 27.8 between one TPU chip and the T4. % src: recomputed: B12.tpu_over_cpu, B11.tpu_over_gpu
Its limitations section said that the Colab values came from \vonequote{the same measurement campaign as the TPU
numbers.}\ Our records do not support this claim.
The T4 and the CPU runtime ran in two Colab runtimes on 8 August 2026, at 04:50 and 05:19 UTC. % src: recomputed: AUG.gpu.first_epoch_utc, AUG.cpu.first_epoch_utc
The baseline run on the TPU did not record its date. Its result was committed at 17:23 UTC on the same day. % src: recomputed: V2.first_commit
The three values also came from different copies of the script (Table~\ref{tab:platforms}).
In September, the T4 took a mean of 6.57~s per steady-state call instead of 13.086~s. % src: recomputed: SEP.gpu.0915.mean, B.gpu.steady
The CPU runtime took means of 348.9~s to 359.3~s instead of 212.113~s. % src: recomputed: V2.sep_cpu_steady_min, V2.sep_cpu_steady_max, B.cpu.steady
This version gives only the smallest ratio over all sessions (Section~\ref{sec:res-platforms}).
In every session, the mean steady-state call on the T4 took more than 13.9 times as long as the recorded call on one TPU chip. % src: recomputed: V2.min_tpu_over_t4_with_max_tpu
On the CPU runtime, it took more than 450 times as long. % src: recomputed: V2.min_tpu_over_cpu_with_max_tpu
Version 1 also gave a ratio between the T4 and the CPU runtime. We removed it.

\subsection*{Speedup on Eight Chips}
Version 1 reported that multi-query \texttt{pmap} on eight chips reached 6.92 times the throughput of the baseline run on the TPU, with a parallel efficiency of 86.5\%. % src: recomputed: MQ.speedup_file, MQ.efficiency_pct
The two runs differed in input and in timed region. This version takes the speedup only from the matched grid. There,
every point used the same script, input family, and timed region.
Eight chips gave 6.53 to 7.91 times the throughput of one chip. % src: recomputed: V2.speedup_8over1_min, V2.speedup_8over1_max
Version 1 explained the lower parallel efficiency at 100 residues by a fixed overhead per call. This version gives the extra time per call and leaves its cause open (Section~\ref{sec:res-pmap}). % src: results/sweep/scaling_law_data.json -> grid[0].length
Version 1 also wrote that \vonequote{distinct per-chip HBM allocations confirm the work is distributed.}\ Memory shows
only that every chip held data. The division of the work follows from \texttt{jax.pmap}, which runs one replica per
device~\cite{jaxpmap}.

\subsection*{First Call and Profiler}
Every August version of the single-query script that we have runs the first call with the profiler. The September
version adds an option that turns the profiler off. The TPU runs did not record the duration of the profiler teardown.
In September, the profiler teardown took 31.7~s to 31.9~s on the CPU runtime and 23.7~s on the T4. % src: recomputed: V2.sep_cpu_tail_min, V2.sep_cpu_tail_max, SEP.gpu.0915.profiler_tail_s
Version 1 gave the TPU first call as 59.1 times the steady-state call, a ratio that it called uncorrected. % src: recomputed: B15.tpu_first_over_steady
This version gives the first call as 27.78~s and gives no ratio. % src: recomputed: B.tpu.first
Table~\ref{tab:platforms} marks that this call probably includes the profiler teardown.
Version 1 also read the compilation-cache experiment as showing that more than half of the first call survives the
cache. Both first calls of that experiment include the profiler teardown, so this version does not attribute the 15.19~s of the second process to compilation. % src: results/sweep/compilation_cache.json -> warm.first_predict_seconds
For the T4 and the CPU runtime, version 1 subtracted the profiler teardown from the August first calls. This version
uses the September runs without the profiler instead.

\subsection*{Profiler Trace}
Version 1 reported a trace of one TPU first call.
In that trace, a JAX function that traces and compiles had 12.55~s of self time out of 16.56~s. This share is about 76\%. % src: profiling/trace_analysis.md ("16.56s", "12.55s"); recomputed: TR.share_pct
We did not keep the trace file, and only a written analysis of it remains. No reader can check the value, so we
removed it.

\subsection*{Automatic Partitioning}
Version 1 named this experiment \vonequote{automatic sharding} and discussed it in terms of GSPMD. It cited the GSPMD
paper for the claim that a partitioner without annotations falls back to replication. That paper does not state this.
The runs did not record which partitioner ran. This version reports the experiment as a run without sharding
annotations. It does not say whether the call ran on one chip or on every chip (Section~\ref{sec:res-auto}).

\subsection*{Cost}
Version 1 found that the slice with one chip working cost about as much per query as the T4.
The values were \$1.25 and \$1.27 per 1000 queries. % src: recomputed: COST.pod.per_1k, COST.gpu.per_1k ("per 1000 queries")
This parity held only for the T4 time of August. The T4 price came from a survey that version 1 did not cite. The CPU
price was our estimate for a machine with four virtual CPUs. The September CPU runtime had two, and the August sessions
did not record theirs.
Version 1 also priced hypothetical slices of one, two, and four chips in a table. This version prices only the TPU, at Google's list price of \$1.20 per chip-hour. % src: recomputed: V2.price_chip_hour
It gives one conditional example for a chip rented alone (Section~\ref{sec:res-cost}).

\subsection*{\AFthree{}}
Version 1 had a section and an appendix on \AFthree{}. The release we tested did not accept the TPU as a platform.
These runs do not address the question of this paper, which concerns the TPU slice. We removed the section, the
appendix, and all their numbers.

\subsection*{Other Corrections}
Version 1 stated that the sequence-length sweep behind the power-law fit covered 100 to 1000 residues. The fit uses the
matched grid.
The sequence-length sweep covered 60 to 500 residues. % src: results/sweep/scaling_law_data.json -> grid[*].length; results/sweep/sequence_length_sweep.json -> [*].num_residues
Version 1 wrote that none of the cited work reported TPU experiments, but it cited a model trained on TPUs. Version 1
described a controlled comparison with identical code on the three platforms. AlphaFold2's source ran unmodified, but
the scripts differed (Table~\ref{tab:platforms}). Version 1 called the copy of the script in the August notebooks
uncommitted.
The notebooks in commit 2755250 hold it. % src: recomputed: V2.first_commit
Seven works cited in version 1 do not make the statements that version 1 attributed to them. We corrected these statements or
removed the citations. Version 1 treated the files in \texttt{results/sweep/} as raw results. They are summaries of the
Job output (Section~\ref{sec:records}). The abstract now states that the weights were random. We also renamed three
terms. A \vonequote{pod} is now a slice, a \vonequote{protein} is a query, and \vonequote{automatic sharding} is
automatic partitioning.
